Signal averaging is the de facto method for improving the signal-to-noise ratio (SNR) in laser ultrasonic testing, yet it imposes fundamental constraints on both inspection throughput and material integrity. Because laser-generated ultrasound is inherently low-energy, the detected signals are typically dominated by background noise, and the SNR improves only as the square root of the number of averages \cite{liu2022deep}. Doubling the averaging count from 256 to 1024, for instance, yields only approximately 3 dB of SNR gain, while doubling the acquisition time and doubling the cumulative laser fluence incident on the specimen surface. This presents a severe trade-off: achieving SNR levels adequate for reliable defect detection---typically requiring 256 to 1024 averages in practice---demands inspection times that render real-time or high-throughput screening impractical, and laser fluence levels that risk surface micro-ablation or coating damage on sensitive components \cite{yang2023unsupervised}. The problem is further compounded for in-situ inspection scenarios, where access is limited, the component surface may be fragile or previously machined, and repeated laser exposure is either physically undesirable or contractually prohibited.

Conventional denoising methods offer limited relief. Linear filters such as Wiener filtering and wavelet-based thresholding can suppress noise but frequently distort diagnostically relevant acoustic features---arrival times, peak amplitudes, and waveform morphology---upon which defect detection depends \cite{chen2019wavelet}. More advanced patch-based methods like BM3D achieve strong noise suppression in image domains but struggle to preserve the sparse, transient character of ultrasonic waveforms in the time domain \cite{dabov2007bm3d}. Existing deep learning approaches show promise for general denoising \cite{zhang2017denoising} but are typically trained and evaluated using data from extensively averaged reference signals, creating a performance gap when applied to the low-averaging conditions encountered in practical inspection \cite{liu2022deep}. The result is a critical gap: inspection scenarios that demand high throughput, involve sensitive or coated surfaces, or require in-situ evaluation cannot rely on extensive averaging, yet current denoising alternatives fail to provide the necessary SNR improvement without sacrificing the acoustic features essential for defect characterization. There is therefore a pressing need for denoising methods capable of delivering reliable signal quality under substantially reduced averaging, enabling practical laser ultrasonic inspection without the surface damage and throughput penalties of conventional signal averaging.
